\section{Introduction}
\citet{lucca2015} document large average excess equity returns in the 24 hours before scheduled FOMC policy statements, about 49~bps in 1994--2011, and subsequent work ties the premium to uncertainty \citep{hu2022,martello2018}. Three questions follow. Does the premium survive publication and the shift to press conferences at every meeting after 2019? Is it concentrated in high-uncertainty meetings, as the uncertainty-pricing account of \citet{hu2022} predicts? And are the original findings robust to how the announcement time is measured? A premium that shrinks after publication would be consistent with the decay of published predictors \citep{mclean2016}.

We fix the specification, state threshold and inference procedure in a pre-analysis plan, estimate on 2010--2019 intraday futures data, and run the 2020--2026 holdout once. In the holdout the mean excess window return is 38.3~bps and exceeds that on matched-VIX non-meeting days by 55.5~bps ($p=0.002$). The anxious-state premium is 37.8~bps but is not distinguishable from zero with year-clustered wild-bootstrap inference ($p=0.149$), and the anxious-minus-calm contrast is $-2.6$~bps ($p=0.946$), against $+19.2$~bps in development. With seven clusters and nine calm meetings the holdout has limited power, so we report bounds rather than evidence of absence.

A second contribution is measurement. Checking each development meeting against Federal Reserve release pages, we find that eight 2011--2012 meetings with press conferences released the statement at 12:20--12:35 ET rather than 14:15 ET, so the frozen window end followed the announcement for those meetings. We correct the development benchmark in a documented post-freeze amendment. The corrected development mean is 2.8~bps rather than 6.2, and the 2010--2014 versus 2015--2019 decline disappears. The holdout, whose statements are released at 14:00 ET, is unaffected. The paper is a disciplined out-of-sample check, not a causal estimate.